\documentclass[10pt,a4paper]{amsart}
\usepackage[margin=19mm]{geometry}
\usepackage{amsmath,amssymb,mathtools}
\usepackage[scr=rsfs,cal=euler]{mathalfa}
\usepackage{xfrac}
\usepackage[unicode,hidelinks,bookmarks=false]{hyperref}
\newcommand{\ie}{i.e.\ }
\newcommand{\cf}{cf.\ }
\newcommand{\resp}{respectively\ }
\newcommand{\Subsection}{\subsection}
\providecommand{\nobreakdash}{\nobreak\hbox{-}}
\setlength{\emergencystretch}{2em}
\multlinegap0pt
\usepackage{amssymb}
\usepackage{xcolor}
\usepackage{enumerate}
\newtheorem{theorem}{Theorem}
\newtheorem{corollary}{Corollary}
\title{Sharp estimates for the Fourier transform of surface-carried measures and maximal operators associated with hypersurfaces in $\mathbb{R}^4$ with vanishing Gaussian curvature.}
\author{Isroil A.\, Ikromov, \, \, \, Gayrat Toshpulatov}
\date{Source: arXiv:2602.18163v2 (CC BY 4.0)}
\begin{document}
\maketitle

\noindent\emph{Source and licence.} Sharp estimates for the Fourier transform of surface-carried measures and maximal operators associated with hypersurfaces in $\mathbb{R}^4$ with vanishing Gaussian curvature.. Version arXiv:2602.18163v2 (CC BY 4.0), distributed by arXiv under the Creative Commons Attribution 4.0 International licence, http://creativecommons.org/licenses/by/4.0/, as recorded in the arXiv licence metadata for this exact version and shown on the abstract page https://arxiv.org/abs/2602.18163v2. Full licence notice: \texttt{licenses/arxiv-2602.18163-CC-BY-4.0.txt}.\par\smallskip

\noindent\textit{Editorial scope.} Counted unit: the article's theorem Bon (source0.tex lines 399--407). The article's complete original proof of it is delivered with it (source0.tex lines 553--578, one \texttt{proof} environment of 26 lines, which the article places away from the statement). No mathematics is written, repaired or re-derived: the statement and the proof are the article's own lines, in the article's own notation, and no retained line is substituted or re-flowed.  Retained uncounted as the source-local context the proof needs: lines 410--413; lines 548--551.  Nothing was omitted from any retained span, so the package carries no omission record.  No retained line contains a citation, so the wrapper carries no bibliography and no bibliography entry.  No retained line contains a figure environment, an image-inclusion command or drawing code, so no image is delivered and the package declares no asset dependency.  Editorial formatting only, in addition: an AMS \texttt{amsart} preamble that restates the article's own declarations and macros used by the retained lines, namely the environments \texttt{theorem}, \texttt{corollary} on separate counters, together with the standard packages those lines need.  The retained environments print in this document as Theorem 1, Corollary 1 in order of appearance, whereas the article prints the same items with the numbering of its own full document.  The complete native source of the article as distributed in the arXiv e-print for this exact version is bundled unchanged as \texttt{sources/195016/source0.tex}.
\begin{theorem}\label{th:adap}
Let  $\phi:\mathbb{R}^3\to \mathbb{R}$ be a non-trivial polynomial function such that $\phi(0)=0$ and $\nabla \phi(0)=0.$  Assume $$\det{(D^2\phi(x))}=0$$
holds  for all $x=(x_1,x_2,x_3)\in \mathbb{R}^3$.   Then there exists an adapted coordinate system to $\phi$ in $\mathbb{R}^3.$ Moreover, let  $A\in \mathbb{R}^{3\times 3}$ be the matrix given in  Theorem \ref{Bon}. If $\phi(Ax)$ only depends on one variable, then the given coordinate  system $x$ is adapted. If it depends on two variables, then an adapted coordinate system can be found by either $y_1 = x_1$, $y_2 = x_2 - \psi(x_1)$ or $y_1 = x_2$, $y_2 = x_1 - \psi(x_2)$, where $\psi$ represents a real-valued analytic function. If $\phi(Ax)$ has a form as in \eqref{form}, then an adapted coordinate system can be obtained   by applying a linear transformation of variables.
\end{theorem}

\begin{corollary}\label{esseen}
Let  $\phi:\mathbb{R}^3\to \mathbb{R}$ be a  polynomial function  satisfying the conditions of Theorem   \ref{th:adap}.  Then, after  a
local analytic  change of variables,   $\phi$ depends on at most two variables.
\end{corollary}

\begin{theorem}\label{Bon}
Let  $\phi:\mathbb{R}^3\to \mathbb{R}$ be a  polynomial  such that $\phi(0)=0$ and $\nabla \phi(0)=0.$ Assume  $$\det{(D^2\phi(x))}=0$$
holds   for all $x=(x_1,x_2,x_3)\in \mathbb{R}^3$.  Then
then there exists an  invertible matrix $A\in \mathbb{R}^{3\times 3}$ such that $\phi(Ax)$ either depends on at most two  variable $x_1$ and $x_2$ or has the form
\begin{equation}\label{form}
\phi(Ax)=Q_1(x_1)+Q_2(x_1)x_2+Q_3(x_1)x_3,
\end{equation}
where $Q_1,Q_2,Q_3:\mathbb{R}\to \mathbb{R}$ are polynomial functions.
\end{theorem}

\begin{proof}[{Proof of Corollary \ref{esseen}}] Theorem \ref{Bon} shows that there is an invertible matrix $A$ such that $\phi(Ax)$ depends on at most two variables or has the form in \eqref{form}. If $\phi(Ax)$ depends on at most two variables, then the corollary is proven. It remains  to consider the case  where  $\phi(Ax)$ has the form in  \eqref{form}. As in the proof of Theorem \ref{th:adap} we can consider two cases. In Case 1, we  change the variables
\begin{align*}
y_1&=x_1|\tilde Q_1(x_1)+x_1^{\nu_2-\nu_1}x_2 \tilde{Q}_2(x_1)+x_1^{\nu_3-\nu_1}x_3\tilde{Q}_3(x_1)|^\frac1{\nu_1},\\
y_2&=x_2,\\
y_3&=x_3.
\end{align*}
One can check that this change of variables is analytic in a sufficiently small neighborhood of the origin.
 Then, in this new coordinate system $y,$ $\phi(Ax)$ equals $\mathrm{sign}(\tilde Q_1(0))y_1^{\nu_1}.$ So, it depends on one variable $y_1$.


As shown in the proof of Theorem \ref{th:adap}, in Case 2 there is an invertible matrix $B$ such that
\begin{align*}
\phi(Bz)=&z_1^{\nu_2}z_2+z_1^{\nu_1}(\tilde Q_1(z_1)-c_1)-\frac{c_1}{c_2}z_1^{\nu_2+1}(\tilde Q_2(z_1)-c_2)\\ &+\frac{z_1^{\nu_2}z_2}{c_2}(\tilde Q_2(z_1)-c_2)+\left(z_1^{\nu_3}(\tilde Q_3(z_1)-c_3)-\frac{c_3z_1^{\nu_2} }{c_2}(\tilde Q_2(z_1)-c_2)\right)z_3.
\end{align*}
 It is sufficient to use the following change of variables
 \begin{align*}
 y_1&=z_1,\\
  y_2&= z_2+z_1^{\nu_1-\nu_2}(\tilde Q_1(z_1)-c_1)-\frac{c_1}{c_2}z_1^{\nu_2+1-\nu_1}(\tilde Q_2(z_1)-c_2) \\ &\quad +\frac{z_2}{c_2}(\tilde Q_2(z_1)-c_2)+\left(z_1^{\nu_3-\nu_2}(\tilde Q_3(z_1)-c_3)-\frac{c_3 }{c_2}(\tilde Q_2(z_1)-c_2)\right)z_3,\\
  y_3&=z_3.
 \end{align*}
Obviously, this change of variables is  analytic  in a sufficiently small neighborhood of the origin. In this new coordinate system $y$, $\phi(Bz)$ equals $ y_1^{\nu_2} y_2.$
So, it depends on two variable $y_1$ and $y_2$.


%The proof in Case 3 follows similarly.
  \end{proof}

\end{document}
